\subsection{Implications to AI Safety}

While iterative deployment shares links with RL fine-tuning, it also brings
new concerns about AI safety. Unlike standard RL training, where reward functions can be explicitly
designed to encode human preferences and safety constraints, iterative
deployment of models in the wild relies on implicit signals from post-deployment usage. Curation done indirectly through user interactions with previously deployed models behaves as an opaque reward
function which can be difficult to understand and control.
In RL
for alignment, reward functions are often used to align the model with safety
constraints and goal specifications. In iterative model deployment, the implicit reward
functions could clash with the explicitly specified rewards in
model alignment. This could lead to unexpected problems in future model
generations, as the indirect curation used might lead to large gradient steps in the opposite direction to the gradients induced by the safety training, for example.
Overall, this raises new alignment challenges.


Another concern is bias during validation. If the external validator (e.g., a tool,
or a human curator) has unintended or malicious biases, these biases might
accumulate over the generations. Later models might then optimize for harmful
properties that diverge from the original goals and also from safety
constraints.

A final property to note is model collapse \citep{shumailov-et-al-nature2024}. Data
curation might delay model collapse by filtering for valid traces, it is still
unknown whether this fully prevents collapse. If collapse
occurs, important capabilities can degrade in ways that are difficult to
detect until deployment. Note, however, that it is also unclear whether
longer training periods using RL techniques, such as PPO or GRPO, lead to model
collapse or not.
